% ======================================================================
% boards/rpi4.tex — Raspberry Pi 4 (BCM2711, Cortex-A72) のマルチコア化
% 実測: 2026-07-19, HTTP /bench (192.168.3.100), production kernel md5 4cf1a50a
% ======================================================================

Raspberry Pi 4（BCM2711、Cortex-A72 $\times$ 4、AArch64）へ、第 \ref{sec:design} 節と
同一のワーカープール型 SMP を移植した。共通設計・測定方法はボードに依存しないので、
ここでは\textbf{Pi 4 固有の差分と実測値}のみを述べる。

\subsection{対象と前提}

BCM2711 は 4 つの Cortex-A72 を持つ。Pi 3（A53、32bit）と異なり\textbf{AArch64
（64bit）}で動作する。ファームウェアはコアを EL2 に落とすため \texttt{boot.S} で
EL2$\to$EL1 に遷移させる。D-cache は第 \ref{sec:design} 節の前提どおり\textbf{OFF}
（\texttt{SCTLR.C=0}、実測 \texttt{SCTLR}=\texttt{0x30d01801}）であり、ロックフリー
調停が RAM 直行で成立する。この D-cache を敢えて ON にする独立実験は第
\ref{sec:dcache-a72} 節に示した。

\subsection{コア起動：spin table}
\label{sec:rpi4-bringup}

Pi 4 のファームウェア（armstub8）はコア 1--3 を物理アドレス
\texttt{0xe0}/\texttt{0xe8}/\texttt{0xf0} でスピンさせて保持する。各スロットへ
エントリ関数アドレスを書き \texttt{SEV} で起こす（Pi 3 の mailbox、Pi 5 の PSCI と
対応する差分）。起床した各コアは \texttt{mmu\_enable\_secondary} で MMU＋I-cache を
コア 0 と同一構成にし、\texttt{cores\_online = 4} を実機で確認した。

\subsection{実測 1：計算律速はスケール、ただし配分方式に強く依存}

\begin{table}[htbp]
\centering
\caption{Pi 4：計算律速のコア数スイープ（\route{/bench}、\texttt{cores=1..4}、全て \texttt{agree=yes}）}
\label{tab:rpi4-compute}
\small
\begin{tabular}{@{}lccccl@{}}
\toprule
ベンチ & 1 & 2 & 3 & 4 & 備考 \\
\midrule
\bench{dining n=5}      & \spd{0.99} & \spd{1.99} & \spd{3.00} & \spd{3.99} & 完全に線形（理想） \\
\bench{primes n=300000} & \spd{0.99} & \spd{1.61} & \spd{2.30} & \spd{3.00} & 連続分割で概ね均衡 \\
\bench{nqueens n=12}    & \spd{0.99} & \spd{1.60} & \spd{1.79} & \spd{1.85} & 静的分割では頭打ち \\
\bottomrule
\end{tabular}
\end{table}

\noindent
\bench{dining}（独立テーブルの均一コスト）は Pi 3・Pi 5 と同じく理想の
\spd{3.99} を示す。\bench{nqueens} は 4 コアで \spd{1.85}——だがこれは
分割方式（連続 vs インターリーブ）の問題ではなく\textbf{静的分割そのものの限界}で
あることを、第 \ref{sec:cross-fair} 項の実験（連続・第 1 行・第 2 行の三方式が
すべて \spd{2.0} 前後）で確認した。動的負荷分散を要する Level 3 の対象である。

\subsection{実測 2：メモリ律速はスケールしない}

\begin{table}[htbp]
\centering
\caption{Pi 4：メモリ書込律速（\bench{fill}、8 パス、\texttt{cores=4}）}
\label{tab:rpi4-fill}
\small
\begin{tabular}{@{}rrrl@{}}
\toprule
語数 & 1 コア \textmu s & 4 コア \textmu s & 速度向上 \\
\midrule
64      & 3     & 7     & \spd{0.42} \\
256     & 8     & 12    & \spd{0.66} \\
1{,}024 & 28    & 35    & \spd{0.80} \\
4{,}096 & 110   & 100   & \spd{1.10} \\
16{,}384 & 438  & 307   & \spd{1.42} \\
65{,}536 & 1{,}749 & 1{,}269 & \spd{1.37} \\
262{,}144 (=1\,MB) & 6{,}997 & 5{,}050 & \spd{1.38} \\
\bottomrule
\end{tabular}
\end{table}

\noindent
Pi 3 と同じ二つの現象が現れる：\textbf{大きな仕事でも \spd{1.4} で頭打ち}（1\,MB で
\spd{1.38}）、\textbf{小さな仕事は並列化すると遅くなる}（64 語で \spd{0.42}）。
D-cache OFF ゆえメモリ書込が RAM 帯域律速になるためで、計算律速の \spd{3.99} とは
比較にならない。

\subsection{実測 3：プールの往復コスト}

空ジョブ（\bench{null}）の 4 コア投入は実測\textbf{約 5\,\textmu s}である（1 コアは
0\,\textmu s）。この往復コストを下回る仕事は並列化すると損をする（表
\ref{tab:rpi4-fill} の 64 語 \spd{0.42} が実例）。Pi 3 の 7--11\,\textmu s より
小さいのは A72 のクロック・キャッシュがコア間シグナリングを速めたためである。

\subsection{小括}

Pi 4（A72）は、\textbf{計算律速は理想的にスケール（\spd{3.99}）・メモリ律速は
スケールしない（\spd{1.38}）}という Pi 3 の結論をそのまま再現した。相違は絶対速度
（A72 は A53 より速い）と、プール往復の小ささ、そして AArch64／spin table という
起動機構のみである。「並列化すべき対象を分ける境界」は Pi 3 と同一であった。
